\section{心理学と計算機}
\label{s01_Introduction}
\subsection{授業内容}
\subsubsection{概要}
心理統計の授業では，多くの数値を一気に扱うため計算機を用いる。Excel/Numbers/Calcsのような表計算ソフトでは不十分なため，Rと呼ばれる統計専門のソフトウェアを用いる。Rの使い方については次回以降に詳細が論じられるが，それに先立って計算機の基礎的な概念，操作方法を理解する。
\subsubsection{コマ主題細目}
\begin{description}

	\item[電子計算機の基礎] 計算機には5つの機能がある。すなわち，入力，出力，演算，制御，記憶装置とよばれるもので，対応するデバイスの組み合わせで計算機が構成される。また，BIOS，OS，アプリケーションの違いなど，計算機の内部構造についての理解は，人間をコンピュータとみなして考えるとき，メタファとして有用である。
	
	\item[情報の単位] 電子計算機はすべて0/1の情報として扱う。さまざまな単位と，対応する装置ではどの程度の大きさを扱えるのかについて，基本的な知識を身につける。ビット，バイト，キロバイトからテラバイトまでの関係性や，テキストデータと画像・動画データの容量の違いなどを理解することで，扱うデータの規模感を把握する。

	\item[ファイルの種類と拡張子] コンピュータで扱うファイルには様々な種類があり，拡張子によって区別される。特に統計解析で扱うことの多いテキストファイル(.txt)，CSV形式(.csv)，Excelファイル(.xlsx)の違いを理解し，それぞれの特徴と用途を把握する。また，Rで扱うデータファイル(.RData)やスクリプトファイル(.R)についても概説する。

	\item[ファイルの場所] コンピュータ内でのファイル管理の仕組みとして，ディレクトリ(フォルダ)階層構造について学ぶ。カレントディレクトリ，相対パス，絶対パスなどの概念を理解し，Rを使った解析においてファイルの読み込みや保存をする際に必要となる基礎知識を習得する。クラウドストレージの利用についても触れ，バックアップの重要性を理解する。

	\item[キーボードとショートカット] 効率的なコンピュータ操作のために，キーボードの配置やタイピングの基本，よく使うショートカットキー(コピー・ペースト，検索，保存など)について学ぶ。特にRやRStudioでの作業に役立つキーボード操作を中心に解説し，実際に操作しながら習得していく。


\end{description}
\subsubsection{キーワード}
\begin{itemize}
	\item コンピュータの基本構造(入力・出力・制御・演算・記憶)
	\item ファイル形式とデータ管理
	\item 効率的なPC操作
\end{itemize}
\subsection{授業情報}
\paragraph{コマの展開方法}
講義
\paragraph{標準シラバスにおける位置づけ}
該当しません
\subsubsection{予習・復習課題}
\paragraph{予習}
大学のPCルームやメディアセンターで利用できるコンピュータの基本操作方法を確認しておくこと。自分のノートPCを授業に持参する場合は，事前にOSのバージョン，ストレージの空き容量，メモリ容量を確認しておくこと。また，普段使用しているデバイス(スマートフォンやタブレット)と比較して，PCではどのような操作が異なるかについて考えてみると良い。
\paragraph{復習}
授業で紹介したファイル操作(作成・保存・移動・削除)を自分のPCで実践してみること。特に，新しいフォルダを作成し，その中にテキストファイルを作成・保存する練習をすること。また，キーボードショートカット(Ctrl+C, Ctrl+V, Ctrl+Z, Ctrl+S など)を意識的に使用する習慣をつけるよう心がけること。次回の授業までに，Rとそのインターフェースとなるソフトウェア(RStudio)をインストールしておくこと。
